\documentclass{article}

\usepackage{amsmath}

\title{Ch. 26 HW \#1}
\author{Principles of Macroeconomics}
\date{\today}

\begin{document}

\maketitle

\section*{Quick Quiz}

\begin{enumerate}

\item (b) -- The U.S. has grown at roughly 1.70\% per year. To figure out roughly how long real GDP would take to double, we can solve
	\begin{gather*}
	1.02^x (\text{GDP}) = 2(\text{GDP}) \\
	1.02^x = 2 \\
	x\log{1.02} = \log{2} \\
	x = \frac{\log{2}}{\log{1.02}}
	x = 35
	\end{gather*}

\item (c) -- U.S. real GDP per capita was roughly \$64,000 while Indian real GDP per capita was roughly \$6,400.

\item (a) -- China grew at roughly 2.56\% while the U.K. grew at 1.40\%

\item (b) -- Human capital boosts productivity which in turn boosts real GDP per capita -- average real incomes.

\item (d) -- See the case study in 26.2.b.

\end{enumerate}

\section*{Questions for Review}

\begin{enumerate}

\item The level of a nation's GDP measures its total income, total expenditure, and total value-added simultaneously. The growth rate of GDP measures how quickly those three quantities are increasing. In the short term, I would rather live in a nation with a high level of GDP because that nation will have a higher average income, and therefore it's likely to have a higher standard of living. In the long term, I would rather live in a nation with a high growth rate because that nation will eventually have a higher average income.

\item \phantom{x}

	\begin{itemize}
	
	\item Physical capital per worker -- the amount of equipment and structures used to produce goods and services relative to the number of workers
	
	\item Human capital per worker -- the knowledge and skills that workers have
	
	\item Technological knowledge -- society's understanding of how to best produce goods and services
	
	\item Natural resources per worker -- the environment assets that an economy has relative to the number of workers
	
	\end{itemize}

\item A college degree is a form of human capital because it represents a level of knowledge and skill that a worker has obtained. 

\end{enumerate}

\end{document}